\section{Datasets and Splits}
\label{sec:datasets_and_splits}

\subsection{Datasets}
\label{sec:datasets}

The pipeline developed in this work is based on four datasets, each with a specific role in the experimental protocol: (i) LOL-v2~\cite{DBLP:journals/tip/YangWHWL21}, (ii) LOL-v1~\cite{DBLP:conf/bmvc/WeiWY018}, (iii) MILLs~\cite{DBLP:journals/corr/abs-2511-15496}, and (iv) ExDark~\cite{DBLP:journals/cviu/LohC19}. All models are trained on a single source domain and are then evaluated on additional datasets without fine-tuning. This choice is intended to measure cross-dataset generalization, i.e., the ability of each model to enhance low-light images whose distribution differs from the training data.

The main training dataset is LOL-v2. This dataset contains both a real-captured subset, hereafter referred to as LOL-v2 Real, and a synthetic subset, referred to as LOL-v2 Synthetic. In the main protocol, only LOL-v2 Real is used, meaning that the low-light/normal-light pairs are acquired from real scenes rather than obtained through synthetic low-light degradation. LOL-v2 is a paired dataset: each low-light image is associated with a corresponding normally exposed reference image, which can be used as the target for supervised training and evaluation.

LOL-v1 is used only as a paired cross-dataset evaluation set. It represents a relatively close cross-dataset scenario, because it belongs to the same general benchmark family as LOL-v2.

MILLs is used as a second paired cross-dataset evaluation set. In particular, the low-resolution DSLR split, MILLs-DSLR, is adopted to keep the evaluation computationally manageable. Unlike LOL-v1, this dataset introduces a controlled illumination setting. Its input images are associated with different illumination levels, and each input has a corresponding well-illuminated reference image. This makes it possible to evaluate the models not only with global results over the whole dataset, but also by studying how performance changes as the input illumination becomes more severe.

Finally, ExDark is used as an unpaired stress-test dataset. Here, unpaired means that the dataset provides low-light images without a corresponding normally exposed reference image. Therefore, ExDark is not used for reference-based comparison, but for no-reference evaluation and qualitative failure analysis on real-world dark scenes. Its category-based organization is preserved, meaning that images remain grouped according to the object category provided by the dataset, such as cars, people, animals, or indoor objects. This organization may help inspect whether some failure cases are related to the type of scene or object.

\subsection{Dataset Licenses and Terms of Use}

Since this project relies on publicly available datasets, their licenses and terms of use were reviewed to ensure that each dataset is used consistently with the conditions specified by its original authors.

For LOL-v1\footnote{\href{https://daooshee.github.io/BMVC2018website/}{Official Source for LOL-v1}} and LOL-v2\footnote{\href{https://github.com/flyywh/CVPR-2020-Semi-Low-Light}{Official Source for LOL-v2}}, no explicit dataset license or terms of use were found in the official sources consulted and linked in this paper.

MILLs\footnote{\href{https://huggingface.co/datasets/davidserra9/MILL}{Official Source for MILLs}} is distributed through its Hugging Face dataset page, where the dataset metadata reports the license as MIT. This permissive license allows use, modification, and redistribution, provided that the copyright notice and permission notice are included in copies or substantial portions of the material.

ExDark\footnote{\href{https://github.com/cs-chan/Exclusively-Dark-Image-Dataset}{Official Source for ExDark}} is subject to dataset-specific usage restrictions. Although the repository reports the project as open source under the BSD-3 license, it also states that commercial usage requires contacting the authors. More specifically, the dataset \texttt{README} indicates that the ExDark dataset may be used only for non-commercial research purposes.

All datasets are used exclusively within the scope of this non-commercial academic work. Any applicable license terms or usage restrictions are respected, and the corresponding dataset papers are cited in the bibliography.

To make these conditions explicit, a dedicated dataset notice is provided both in the project repository and in the accompanying release. The included datasets remain subject to the original authors' licenses and terms of use, and are provided solely to support non-commercial academic reproducibility.

\subsection{Split Construction}

Split construction follows the role assigned to each dataset, as discussed in Section~\ref{sec:datasets}. For the main source domain, LOL-v2 Real, the official training split is used to derive the training and validation partitions. In particular, a fixed fraction of the official training samples is deterministically held out as validation data using a fixed random seed. The official test split is kept untouched for the final in-domain evaluation. In this context, ``in-domain'' means that training, validation, and testing images come from the same dataset and acquisition setting.

The synthetic subset of LOL-v2 is kept separate from the main protocol and is used only for an additional ablation study, aimed at evaluating whether synthetically degraded low-light examples improve or degrade generalization. When this subset is used, a validation partition is deterministically derived from the official training split using the same procedure adopted for LOL-v2 Real, ensuring consistency in the split construction process.

For the remaining datasets, no additional train or validation partitions are created. Instead, all available samples are assigned to a single evaluation partition. For datasets that already provide predefined splits, such as LOL-v1 and MILLs, the original partitions are merged into a unified test set. ExDark is likewise converted into a single test split, while preserving the original category labels as metadata.

Following the split construction procedure, the resulting split sizes are summarized in Table~\ref{tab:dataset_splits}.

\begin{table}[t]
	\centering
	\caption{Dataset split sizes used in the pipeline. Counts refer to samples, i.e., low-light input images. For paired datasets, each input sample is associated with a corresponding target image.}
	\label{tab:dataset_splits}
	\begin{tabular}{lccc}
		\toprule
		Dataset & Train & Validation & Test \\
		\midrule
		ExDark & -- & -- & 7363 \\
		LOL-v1 & -- & -- & 500 \\
		LOL-v2 Real & 620 & 69 & 100 \\
		LOL-v2 Synthetic & 810 & 90 & 100 \\
		MILLs & -- & -- & 499 \\
		\bottomrule
	\end{tabular}
\end{table}